% Options for packages loaded elsewhere
\PassOptionsToPackage{unicode}{hyperref}
\PassOptionsToPackage{hyphens}{url}
\documentclass[
  ignorenonframetext,
]{beamer}
\newif\ifbibliography
\usepackage{pgfpages}
\setbeamertemplate{caption}[numbered]
\setbeamertemplate{caption label separator}{: }
\setbeamercolor{caption name}{fg=normal text.fg}
\beamertemplatenavigationsymbolsempty
% remove section numbering
\setbeamertemplate{part page}{
  \centering
  \begin{beamercolorbox}[sep=16pt,center]{part title}
    \usebeamerfont{part title}\insertpart\par
  \end{beamercolorbox}
}
\setbeamertemplate{section page}{
  \centering
  \begin{beamercolorbox}[sep=12pt,center]{section title}
    \usebeamerfont{section title}\insertsection\par
  \end{beamercolorbox}
}
\setbeamertemplate{subsection page}{
  \centering
  \begin{beamercolorbox}[sep=8pt,center]{subsection title}
    \usebeamerfont{subsection title}\insertsubsection\par
  \end{beamercolorbox}
}
% Prevent slide breaks in the middle of a paragraph
\widowpenalties 1 10000
\raggedbottom
\AtBeginPart{
  \frame{\partpage}
}
\AtBeginSection{
  \ifbibliography
  \else
    \frame{\sectionpage}
  \fi
}
\AtBeginSubsection{
  \frame{\subsectionpage}
}
\usepackage{iftex}
\ifPDFTeX
  \usepackage[T1]{fontenc}
  \usepackage[utf8]{inputenc}
  \usepackage{textcomp} % provide euro and other symbols
\else % if luatex or xetex
  \usepackage{unicode-math} % this also loads fontspec
  \defaultfontfeatures{Scale=MatchLowercase}
  \defaultfontfeatures[\rmfamily]{Ligatures=TeX,Scale=1}
\fi
\usepackage{lmodern}
\ifPDFTeX\else
  % xetex/luatex font selection
\fi
% Use upquote if available, for straight quotes in verbatim environments
\IfFileExists{upquote.sty}{\usepackage{upquote}}{}
\IfFileExists{microtype.sty}{% use microtype if available
  \usepackage[]{microtype}
  \UseMicrotypeSet[protrusion]{basicmath} % disable protrusion for tt fonts
}{}
\makeatletter
\@ifundefined{KOMAClassName}{% if non-KOMA class
  \IfFileExists{parskip.sty}{%
    \usepackage{parskip}
  }{% else
    \setlength{\parindent}{0pt}
    \setlength{\parskip}{6pt plus 2pt minus 1pt}}
}{% if KOMA class
  \KOMAoptions{parskip=half}}
\makeatother
\usepackage{longtable,booktabs,array}
\usepackage{caption}
\captionsetup[table]{skip=6pt}
\newcounter{none} % for unnumbered tables
\usepackage{calc} % for calculating minipage widths
\usepackage{caption}
% Make caption package work with longtable
\makeatletter
\def\fnum@table{\tablename~\thetable}
\makeatother
\setlength{\emergencystretch}{3em} % prevent overfull lines
\providecommand{\tightlist}{%
  \setlength{\itemsep}{0pt}\setlength{\parskip}{0pt}}
% Auto-generated by infrastructure/rendering/_pdf_latex_helpers.py::extract_math_font_preamble.
% Minimal header passed to Pandoc via -H so Beamer slides pick up the same math font
% as the combined-PDF path. Adjust the manuscript's preamble.md to change the font globally.
\usepackage{unicode-math}
\setmathfont{latinmodern-math.otf}

% Slide layout defaults for warning-clean scientific decks.
\usepackage{etoolbox}
\IfFileExists{xurl.sty}{\usepackage{xurl}}{}
\IfFileExists{seqsplit.sty}{\usepackage{seqsplit}}{\newcommand{\seqsplit}[1]{#1}}
\protected\def\breakseq#1{\seqsplit{#1}}
\protected\def\breaktt#1{\begingroup\ttfamily\seqsplit{#1}\endgroup}
\setlength{\emergencystretch}{6em}
\tolerance=5000
\hbadness=10000
\hfuzz=1pt
\setlength{\tabcolsep}{2pt}
\AtBeginEnvironment{longtable}{\tiny\renewcommand{\arraystretch}{0.86}\setlength{\tabcolsep}{1pt}}
\AtBeginEnvironment{tabular}{\tiny\renewcommand{\arraystretch}{0.86}\setlength{\tabcolsep}{1pt}}
\AtBeginEnvironment{equation}{\tiny}
\AtBeginEnvironment{equation*}{\tiny}
\AtBeginEnvironment{align}{\tiny}
\AtBeginEnvironment{align*}{\tiny}
\AtBeginEnvironment{itemize}{\footnotesize}
\AtBeginEnvironment{enumerate}{\footnotesize}
\AtBeginEnvironment{description}{\footnotesize}
\setbeamerfont{caption}{size=\tiny}
\setbeamerfont{caption name}{size=\tiny}
\setbeamerfont{normal text}{size=\small}
\setbeamerfont{frametitle}{size=\small}
\setbeamerfont{section title}{size=\footnotesize}
\setbeamerfont{subsection title}{size=\footnotesize}
\setbeamertemplate{section page}{%
  \centering
  \begin{beamercolorbox}[sep=12pt,center,wd=\paperwidth]{section title}
    \parbox{0.86\paperwidth}{\centering\usebeamerfont{section title}\insertsection\par}
  \end{beamercolorbox}
}
\setbeamertemplate{subsection page}{%
  \centering
  \begin{beamercolorbox}[sep=8pt,center,wd=\paperwidth]{subsection title}
    \parbox{0.86\paperwidth}{\centering\usebeamerfont{subsection title}\insertsubsection\par}
  \end{beamercolorbox}
}
\setlength{\abovecaptionskip}{2pt}
\setlength{\belowcaptionskip}{0pt}

% Natbib and cross-reference fallbacks — slides don't load natbib
% or cleveref, but combined-PDF manuscript prose may emit these
% commands. The fallback renders citations as a bracketed key list
% and cross-references as detokenized labels so slides don't fail on
% undefined control sequences or raw underscores. \providecommand is
% a no-op if packages load later.
\providecommand{\citep}[1]{[#1]}
\providecommand{\citet}[1]{#1}
\providecommand{\citealp}[1]{#1}
\providecommand{\citeauthor}[1]{#1}
\providecommand{\citeyear}[1]{#1}
\providecommand{\cref}[1]{\texttt{\detokenize{#1}}}
\providecommand{\Cref}[1]{\texttt{\detokenize{#1}}}

\newtheorem{proposition}{Proposition}
\newtheorem{hypothesis}{Hypothesis}
\usepackage{bookmark}
\IfFileExists{xurl.sty}{\usepackage{xurl}}{} % add URL line breaks if available
\urlstyle{same}
% fallback for those not using the hyperref driver hyperxmp:
\makeatletter
\@ifundefined{xmpquote}{\newcommand{\xmpquote}[1]{#1}}{}
\makeatother
\hypersetup{
  hidelinks,
  pdfcreator={LaTeX via pandoc}}

\author{\texorpdfstring{}{}}
\date{}

\begin{document}

\section{Introduction: Why Active Inference Needs an Auditable Learning
Path}\label{introduction-why-active-inference-needs-an-auditable-learning-path}

\subsection{Problem: prerequisite opacity, numerical drift, and
evidence-status
confusion}\label{problem-prerequisite-opacity-numerical-drift-and-evidence-status-confusion}

\begin{frame}[allowframebreaks]{Historical lineage and model boundary}
\protect\phantomsection\label{historical-lineage-and-model-boundary}
The Free Energy Principle is often situated in a historically bounded
perception-as-inference context that includes Helmholtz's account of
unconscious inference, Bayesian-brain accounts of neural uncertainty,
and tutorials that make prediction error and variational bounds
teachable (Helmholtz 1867; Knill and Pouget 2004; Bogacz 2017; Sprevak
and Smith 2023). Its deeper background is selective rather than a single
direct descent: ancient accounts of demonstration and ordered
instruction, medieval optical work, early-modern method and probability,
and eighteenth-century accounts of association and education all expose
related questions about evidence, perception, uncertainty, and sequence
(Aristotle 1901; Ibn al-Haytham 1989; Quintilian 95 AD; Bacon 1620;
Descartes 1628; Huygens 1657; Newton 1704; Rousseau 1762). Bernoulli, de
\end{frame}

\begin{frame}[allowframebreaks]{Historical lineage and model boundary}
Moivre, Bayes, and Laplace then anchor more specific histories of
probability and inverse-probability reasoning, while Berkeley, Hume,
Condillac, Hartley, Locke, Comenius, and Kant provide distinct
traditions of perception, association, habit, experience, and structured
cognition (Bernoulli 1713; Moivre 1718; Bayes and Price 1763; Laplace
1774; Berkeley 1709; Hume 1739; Condillac 1754; Hartley 1749; Locke
1690; Comenius 1657; Kant 1781). These sources are historical context
and design analogies, not evidence that modern variational free energy,
Bayesian brain models, or Active Inference already existed in their
period; the fuller selective history appears in
sec.~17.

\end{frame}

\begin{frame}[allowframebreaks]{Historical lineage and model boundary}
The modern FEP frames self-organizing systems in terms of variational
free energy, with review work clarifying both the principle's scope and
the need to avoid oversimplified readings (Friston et al. 2006; Friston
2010; Friston et al. 2023; Zhang and Xu 2024). In the discrete
formulations used here, Active Inference supplies model-based rules for
belief updating and policy evaluation: variational free energy is used
for current inference, while expected free energy evaluates policies
under stated preferences and observation models (Friston et al. 2017; Da
Costa et al. 2020; Parr et al. 2022; Pezzulo et al. 2024). The word
\emph{agent} therefore names a modeled decision-making system; it does
not by itself establish consciousness, intelligence, autonomy, or
empirical adequacy.
\end{frame}

\subsection{Situation awareness: theory, model, software, and education
are different
objects}\label{situation-awareness-theory-model-software-and-education-are-different-objects}

\begin{frame}[fragile,allowframebreaks]{Education and delivery boundary}
\protect\phantomsection\label{education-and-delivery-boundary}
The field is easiest to misread when four layers are collapsed into one
label. At the theory layer, the FEP is a broad explanatory and
mathematical framing whose scope, generality, and distinctive empirical
predictions remain matters of active debate (Friston et al. 2023;
Aguilera et al. 2022; Gershman 2019; Hodson et al. 2024). At the model
layer, a discrete Active-Inference process theory specifies beliefs,
likelihoods, preferences, policies, and an inference or planning scheme;
changing those assumptions changes the model's behavior, even when
authors use the same vocabulary (Friston et al. 2017; Da Costa et al.
2020; Parr et al. 2022; Millidge et al. 2021). A software implementation
then adds a third layer: libraries such as \texttt{pymdp} and this
project's kernels are executable realizations with different goals,
tests, and coverage, not interchangeable evidence for the theory or for
\end{frame}

\begin{frame}[fragile,allowframebreaks]{Education and delivery boundary}
biological adequacy (Heins et al. 2022; Oostrum et al. 2025).

The education layer adds a fourth boundary. Active-learning and
prepared-environment proposals make the framework relevant to
instruction, while knowledge-space and curriculum-aware
cognitive-diagnosis work describe stronger learner-model commitments
that require response data and an explicit account of learner state (Di
Paolo et al. 2024; Doignon and Falmagne 1985; Falmagne and Doignon 2011;
Fu and Fang 2025). Learning analytics can support such future
evaluation, but its use also creates privacy, data-protection, validity,
and human-oversight obligations (Liu and Khalil 2023; Sclater and Bailey
2015). SkillTree is a delivery substrate with its own runtime semantics;
it does not convert a curriculum export into an intelligent tutor or
turn a model demonstration into learner-outcome evidence (SkillTree
2026; National Security Agency 2020).
\end{frame}

\begin{frame}[fragile,allowframebreaks]{Education and delivery boundary}

A modeled Active-Inference agent is therefore a formal decision-making
object with specified states, observations, preferences, and policies;
it is not by itself an intelligent system. An intelligent system would
require a separate capability and evaluation argument, while an adaptive
learner-facing tutor would additionally require learner-state data,
intervention logic, and prospective outcome evidence. Active
Skillference occupies the intersection of these layers as a
provenance-bound curriculum artifact: it makes a selected model family
representable, teachable, and inspectable without claiming to settle the
theory, personalize instruction, or measure learning.

\end{frame}

\begin{frame}[fragile,allowframebreaks]{Education and delivery boundary}
That compact summary hides a long prerequisite chain. To interpret
expected free energy, a learner needs KL divergence; to interpret KL
divergence, entropy and expectation; and to work with those quantities,
probability, logarithms, optimization, and the geometry of
distributions. A learner can encounter apex terms such as ``expected
free energy'' or ``epistemic value'' before the mathematical foundations
that make them derivable. The pedagogical risk is conceptual
miscalibration: a learner may know the words before being able to
compute the quantities or distinguish mathematical identities from
process-theory commitments, empirical hypotheses, and philosophical
scope claims. Active Inference tutorials and expected-free-energy
derivation reviews show why the formal sequence cannot be replaced by
slogans, while education-oriented and associative-learning Active
Inference work frames learning as structured exploration of prepared
\end{frame}

\begin{frame}[fragile,allowframebreaks]{Education and delivery boundary}
material, attention, and prediction-error reduction rather than passive
term recognition (Smith et al. 2022; Millidge et al. 2021; Champion et
al. 2026; Di Paolo et al. 2024; Anokhin et al. 2024).

A second problem is quantitative integrity. Hand-maintained numerical
examples can drift from the equations they illustrate. Without an
executable source of truth, a reader cannot tell whether a displayed
value was computed, rounded, copied, or invented. This is a design risk
that the platform addresses; it is not a prevalence claim about every
tutorial or textbook.

\end{frame}

\begin{frame}[fragile,allowframebreaks]{Education and delivery boundary}
A third problem is scholarship drift. A manuscript can accumulate
citations that are syntactically valid but weakly connected to the
claims they are supposed to support, while newer reviews or critiques
change the evidential status of a phrase without changing its syntax.
That matters here because technical critiques of FEP generality and
recent empirical-status reviews, psychology-and-psychiatry progress
reviews, and distinctive-prediction critiques belong in the evidential
context (Aguilera et al. 2022, 2023; Hodson et al. 2024; Badcock and
Davey 2024; Gershman 2019). This project treats the bibliography as part
of the build contract: every cited source is mapped to a support lane
and role, then checked before figures or learner artifacts are emitted.
\end{frame}

\subsection{Approach: provenance-bound prerequisites delivered through
SkillTree}\label{approach-provenance-bound-prerequisites-delivered-through-skilltree}

\begin{frame}[fragile,allowframebreaks]{Approach: provenance-bound
prerequisites delivered through SkillTree}
Active Skillference addresses the three failure modes with one
construction: a validated directed acyclic graph whose nodes are atomic
\emph{skills} and whose edges are explicit prerequisites. A skill is
reachable only after the graph resolves the dependencies that precede
it. The graph spans 8 strata from mathematical foundations to
applications and is checked to be acyclic, connected from foundations to
applications, and free of dangling prerequisites. This follows the basic
premise of prerequisite-sensitive curriculum design: ordering is part of
the learning object (Bloom 1956; Liang et al. 2018). The result is also
interpretable through knowledge-space and learning-space theory, but
only as an authored closure condition over course content; this release
does not infer an individual's knowledge state or implement an adaptive
diagnostic learning-space engine (Doignon and Falmagne 1985; Falmagne
\end{frame}

\begin{frame}[fragile,allowframebreaks]{Approach: provenance-bound
prerequisites delivered through SkillTree}
and Doignon 2011).

The evaluation is correspondingly bounded by four questions. First, does
the curriculum graph satisfy its declared structural and assessment
invariants? Second, can quantitative examples be regenerated from tested
kernels without unbacked or hand-entered result values? Third, can
citation roles, manuscript claims, figures, and generated artifacts be
audited as one evidence surface? Fourth, does the export preserve
SkillTree's ownership of learning paths, quizzes, progress, and scoring?
These questions evaluate the artifact and its delivery contract; they do
not test whether the curriculum is optimal, whether an Active-Inference
model is empirically true, or whether learners improve.

\end{frame}

\begin{frame}[fragile,allowframebreaks]{Approach: provenance-bound
prerequisites delivered through SkillTree}
Every quantitative claim in a skill's content is produced by a tested
computational kernel and inserted through a \texttt{claim:NAME} token.
The platform refuses to publish a skill whose claims are not backed by
code, and a linter flags any bare result number that slips into prose
outside a claim token or LaTeX math. Thus the worked example that
teaches entropy shows 0.693147 nats for a fair coin and 4.60517 nats for
a one-in-a-hundred event because the kernel computed those values.

\end{frame}

\begin{frame}[fragile,allowframebreaks]{Approach: provenance-bound
prerequisites delivered through SkillTree}
Finally, the curriculum is built for delivery. It targets
\textbf{SkillTree}, an open-source training and gamification platform
whose public documentation describes Project -\textgreater{} Subject
-\textgreater{} Skill composition, learning paths, and quiz-based
knowledge checks (National Security Agency 2020, 2026b, 2026a; SkillTree
2026). Active Skillference maps subjects, skills, prerequisite learning
paths, and quizzes onto that data model so the exported artifact can be
inspected and, when the service is available, handed to a real
progress-tracking environment. The citation is a platform-provenance
claim: it identifies the SkillTree substrate and its government
open-source origin, not an endorsement of this curriculum or a security
assurance for any deployment.

\end{frame}

\begin{frame}[fragile,allowframebreaks]{Approach: provenance-bound
prerequisites delivered through SkillTree}
The delivery layer also belongs to a wider community and open-science
setting. The Active Inference Institute's public website repository
describes a static resource hub with audience-specific pathways for
newcomers, learners, researchers, developers, contributors, partners,
and supporters; public repositories and learning resources; and an
explicit boundary for sanitizing and verifying public projections
(Active Inference Institute 2026). Active Skillference contributes one
interoperable piece to that ecosystem: a provenance-bound curriculum and
SkillTree delivery contract that makes formal examples, source roles,
prerequisites, and release checks inspectable. This is a contribution to
community-facing educational infrastructure, not a claim that the
project represents the Institute, speaks for its contributors, or
replaces its broader open-science programs.

\end{frame}

\begin{frame}[fragile,allowframebreaks]{Approach: provenance-bound
prerequisites delivered through SkillTree}
These design choices turn the paper's contribution from a list of topics
into a set of auditable contracts. Each contribution below is either a
checked structure, a provenance gate, a source-role audit, or a delivery
path into the existing SkillTree runtime.

The contribution map is therefore tied to six evidence classes.
\textbf{Structural evidence} covers graph, prerequisite, assessment, and
artifact invariants. \textbf{Computed evidence} covers values
regenerated by tested kernels and resolved claim tokens.
\textbf{Simulation evidence} covers fixed-input demonstrations whose
behavior is conditional on an explicit model and parameter regime.
\textbf{Literature evidence} covers historical lineage, formal
definitions, pedagogical rationale, platform provenance, and critiques,
each assigned a source role. \textbf{Delivery evidence} covers export,
REST-contract, dashboard, and publication surfaces. \textbf{Scope-limit
evidence} records what the validators and sources cannot establish:
\end{frame}

\begin{frame}[fragile,allowframebreaks]{Approach: provenance-bound
prerequisites delivered through SkillTree}
biological adequacy, learner-state inference, personalization, efficacy,
retention, or transfer. No evidence class is used as a substitute for
another.
\end{frame}

\subsection{Contributions: graph validation, provenance binding,
scholarship audit, and
delivery}\label{contributions-graph-validation-provenance-binding-scholarship-audit-and-delivery}

\begin{frame}[allowframebreaks]{Contributions: graph validation,
provenance binding, scholarship audit, and delivery}
\begin{enumerate}
\tightlist
\item
  \textbf{Structural evidence:} a typed, validated
  \textbf{prerequisite-DAG engine} for curriculum design, including
  cycle detection, topological ordering, depth levels, and connectivity
  checks.
\item
  \textbf{Computed evidence:} a \textbf{content-provenance binding}
  mechanism that ties every learner-facing number to a tested kernel and
  fails the build otherwise.
\item
  \textbf{Curriculum representation:} an implemented \textbf{Active
  Inference curriculum} -- 630 skills across all 8 strata -- authored
  against that engine.
\item
  \textbf{Delivery evidence:} a \textbf{SkillTree exporter and seeder}
  that turns the curriculum into an importable, gamified training
  profile, with live availability treated as an external dependency.
\item
  \textbf{Literature evidence:} a \textbf{scholarship coverage audit}
  that maps manuscript citations to source roles and generates a
  reproducible coverage artifact.
\item
  \textbf{Scope-limit evidence:} an explicit boundary between formal
  model demonstrations, implementation checks, delivery checks, and
  future learner-outcome evaluation.
\end{enumerate}
\end{frame}

\end{document}
